\documentclass[a4paper,12pt]{article}
\usepackage[utf8]{inputenc}
\usepackage[italian]{babel}
\usepackage{amsmath, amssymb}
\usepackage{graphicx}
\usepackage{listings}
\usepackage{xcolor}
\usepackage{hyperref}
\usepackage{bookmark}
\usepackage{titlesec}
\usepackage{fancyhdr}
\usepackage{caption}
\usepackage{geometry}
\usepackage{tcolorbox}
\usepackage{lmodern}


\geometry{margin=2.5cm}
\pagestyle{fancy}
\setlength{\headheight}{16pt}
\fancyhf{}
\renewcommand{\headrulewidth}{1pt}
\renewcommand{\footrulewidth}{0.4pt}
\rhead{\textit{Corso di Object Orientation}}
\lhead{\textbf{Documentazione Progetto Java}}
\cfoot{\thepage}

% Titoli colorati e separatori
\definecolor{sectioncolor}{RGB}{0,70,140}
\definecolor{subsectioncolor}{RGB}{0,120,200}
\titleformat{\section}[block]{\Large\bfseries\color{sectioncolor}\titlerule[0.8pt]}{\thesection}{1em}{}
\titleformat{\subsection}[block]{\large\bfseries\color{subsectioncolor}}{\thesubsection}{1em}{}

% Tema listings moderno
\lstset{
  language=Java,
  breaklines=true,
  basicstyle=\ttfamily\small,
  frame=single,
  backgroundcolor=\color{gray!10},
  keywordstyle=\color{blue!80!black}\bfseries,
  commentstyle=\color{green!50!black}\itshape,
  stringstyle=\color{orange!80!black},
  numbers=left,
  numberstyle=\tiny\color{gray},
  tabsize=2,
  showstringspaces=false,
  captionpos=b
}

% Box decorativi per sezioni chiave
\newtcolorbox{keybox}[1][] {
  colback=blue!5!white,
  colframe=blue!60!black,
  fonttitle=\bfseries,
  coltitle=black,
  left=1mm, right=1mm, top=1mm, bottom=1mm,
  title=#1
}

\hypersetup{
colorlinks=true,
linkcolor=blue,
urlcolor=blue,
pdftitle={Documentazione Progetto Java},
pdfauthor={Gino Pandozzi-Trani}
}

\title{\Huge\textbf{Documentazione del Progetto Java}\\\Large Corso di Object Orientation}
\author{\Large Gino Pandozzi-Trani\\\normalsize Matricola: N86003116}
\date{}


\begin{document}

% Frontespizio elegante
\begin{titlepage}
    \centering
    {\Huge\bfseries Documentazione del Progetto Java\\[1.5ex]}
    {\Large Corso di Object Orientation\\[2ex]}
    \vspace{1.5cm}
    {\large\textbf{Autore:} Gino Pandozzi-Trani}\\[0.5ex]
    {\large\textbf{Matricola:} N86003116}\\[2ex]
    \vspace{1.5cm}
    \rule{0.8\textwidth}{0.8pt}\\[1.5cm]
    {\large\today}
    \vfill
    \begin{flushright}
    \end{flushright}
\end{titlepage}

\thispagestyle{empty}
\newpage
\tableofcontents
\newpage

\section{Contesto e Introduzione}
\begin{tcolorbox}[colback=gray!10, colframe=black!40!black, title=Descrizione generale]
Questo progetto consiste nello sviluppo di un sistema informativo completo per la gestione di un negozio di ortofrutta, realizzato in Java con interfaccia grafica (Swing). Il sistema integra una base di dati relazionale e consente la gestione efficiente di prodotti, clienti, dipendenti e vendite. Le tipologie di prodotto (frutta, verdura, farinacei, latticini, uova, confezionati) sono gestite con attributi specifici: ad esempio, frutta e verdura hanno la data di raccolta, i latticini sia la data di mungitura che quella di produzione, ecc.
\end{tcolorbox}

\subsection*{Funzionalità Implementate}
\begin{keybox}[Funzionalità principali]
\begin{itemize}
    \item Gestione completa dei prodotti e delle relative categorie
    \item Anagrafica clienti con tessera punti e fidelizzazione
    \item Monitoraggio e visualizzazione delle vendite e dei dipendenti
    \item Classifiche dinamiche per vendite e introiti dei dipendenti
    \item Interfaccia grafica intuitiva e moderna (Swing)
\end{itemize}
\end{keybox}

Ogni cliente è registrato e possiede una tessera punti che lo identifica presso il fruttivendolo: per ogni acquisto riceve automaticamente il 10\% del valore della spesa in punti fedeltà. Il sistema consente la ricerca avanzata dei clienti, differenziandoli in base alle categorie di prodotti acquistati e alla quantità di punti ottenuti per ciascuna categoria.

\vspace{0.5em}
\noindent\textbf{Traccia d'esame: Sistema di gestione per un negozio di alimentari}

\begin{tcolorbox}[colback=yellow!10, colframe=orange!80!black, title=Obiettivo della traccia]
Per gruppi da 3: il fruttivendolo ha un certo numero di dipendenti e ad ogni vendita è associato il dipendente che l’ha portata a termine. Il sistema permette di ricercare quale dipendente ha effettuato il maggior numero di vendite in un determinato periodo, distinguendolo da quello che ha portato il maggior introito.
\end{tcolorbox}


\subsection*{Architettura del Sistema}
Il progetto adotta il pattern architetturale \textbf{Model-View-Controller (MVC)}, separando in modo chiaro la logica applicativa, l’interfaccia grafica e l’accesso ai dati. Questa scelta garantisce modularità, manutenibilità e facilità di estensione.

\begin{itemize}
    \item \textbf{Model}: gestisce le entità principali (Prodotto, Cliente, Dipendente, Ordine, Tessera, ecc.) e la logica di dominio.
    \item \textbf{View}: realizza la GUI in Swing, con schermate dedicate all’inserimento, modifica e visualizzazione dei dati.
    \item \textbf{Controller}: gestisce l’interazione tra vista e modello, collegando l’interfaccia grafica alle funzionalità applicative.
\end{itemize}

\section{Login: Admin e Dipendenti}
All’avvio, il sistema presenta una schermata di login dove l’utente inserisce le proprie credenziali. Se l’utente è un amministratore, accede a tutte le funzionalità di gestione e controllo. Se è un dipendente, accede alle funzioni operative e statistiche personali. La sicurezza delle password e la gestione degli errori di autenticazione sono garantite.

\section{Visione degli Elementi}
Il sistema permette la visualizzazione dettagliata di tutti gli elementi gestiti: prodotti, clienti, dipendenti e ordini. Ogni categoria ha una schermata dedicata che mostra le informazioni principali e consente la ricerca e il filtraggio dei dati. L’interfaccia è progettata per essere intuitiva e favorire la consultazione rapida.

\section{Aggiunta e Modifica di Elementi}
L’amministratore e i dipendenti abilitati possono aggiungere nuovi prodotti, clienti e dipendenti tramite apposite schermate di inserimento. È possibile modificare i dati esistenti, aggiornando le informazioni in modo semplice e sicuro. Tutte le operazioni sono tracciate per garantire la correttezza e la trasparenza delle modifiche.

\section{Statistiche e Ricerca Clienti}
Il sistema offre strumenti avanzati per la consultazione delle statistiche di vendita, introiti e performance dei dipendenti. È possibile ricercare i clienti in base a criteri personalizzati, come la quantità di punti, le categorie di prodotti acquistati e la frequenza degli ordini. Le statistiche sono presentate in modo chiaro e sono utili per analisi e decisioni strategiche.

\end{document}
